\documentclass[12pt,a4paper]{article}
\usepackage{iftex}
\ifPDFTeX
  \usepackage[T2A]{fontenc}
  \usepackage[utf8]{inputenc}
\else
  \usepackage{fontspec}
  \setmainfont{Liberation Serif}
  \setmonofont{Liberation Mono}
\fi
\usepackage[russian]{babel}
\usepackage[margin=2cm]{geometry}
\usepackage{amsmath}

\pagestyle{empty}
\setlength{\parindent}{0pt}
\setlength{\parskip}{6pt}

\begin{document}

\begin{center}
{\Large \textbf{Контрольная работа N.2}}\\[4pt]
Вариант 0
\end{center}

\textbf{Время:} задание 1 -- 20 минут; задания 2 и 3 -- по 30 минут.

\hrule

\subsection*{Задание 1 (1 балл)}
Решить уравнение второго порядка, допускающее понижение порядка:
\[
y\,y''+(y')^2=6x.
\]

\subsection*{Задание 2}
Решить уравнение второго порядка методом неопределённых коэффициентов,
проверить найденное решение подстановкой в исходное уравнение и решить
задачу Коши:
\[
y''+2y'+2y=5\sin x, \qquad y(0)=0, \qquad y'(0)=1.
\]

\subsection*{Задание 3}
Решить уравнение второго порядка методом вариации постоянных и проверить
найденное решение подстановкой в исходное уравнение:
\[
y''+y=\tan x, \qquad |x|<\frac{\pi}{2}.
\]

\newpage
\section*{Ответы для самопроверки}

\textbf{1.}\quad
\[
y^2=2x^3+C_1x+C_2.
\]

\textbf{2.}\quad
\[
y=2e^{-x}(\cos x+\sin x)-2\cos x+\sin x.
\]
Для частного решения $y_p=-2\cos x+\sin x$:
\[
y_p''+2y_p'+2y_p=5\sin x,
\]
а также $y(0)=0$, $y'(0)=1$.

\textbf{3.}\quad На интервале $|x|<\pi/2$:
\[
y=C_1\cos x+C_2\sin x-\cos x\ln(\sec x+\tan x).
\]
Для частного решения $y_p=-\cos x\ln(\sec x+\tan x)$:
\[
y_p''+y_p=\tan x.
\]

\end{document}
